\textbf{A second correction:} v1 also claimed "almost everything" determining product quality lives in the tooling vendor's wrapping rather than the model. That's an overgeneralization the surrounding argument doesn't actually need, and it's inconsistent with the model-vendor-supply-chain point made two paragraphs later. The more defensible claim: \textbf{product performance is jointly determined by the model, the data it's given, the harness built around it, the controls in place, and the operating context --- and which of these dominates varies by task.} For some tasks, model choice swamps everything else; for others, orchestration and retrieval quality matter more than which frontier model sits underneath. The reframe's actual point survives without the overclaim: \textbf{an AI tool is not "an AI."} It is software that wraps a model, and the wrapping is a real, evaluable, often-decisive engineering artifact in its own right --- not that the wrapping is \emph{always} the dominant factor.

Once this is visible, "we've built an autonomous agent" often resolves into something much more mundane: "we've built a loop that calls a model and does something with the results." That's not nothing --- the loop can be excellent or terrible engineering --- but it's an ordinary artifact that can be evaluated on its own merits.

\textbf{Guardrail on the reframe:} the model is a replaceable dependency in principle but rarely operationally or legally fungible in a given deployment. Which model vendor and model a tooling vendor has chosen affects accuracy, context capacity, data terms and residency, version stability under silent updates, indemnities, and concentration risk. Evaluate the tooling vendor's wrapping and the model-vendor supply chain behind it as \textbf{two separate questions}.

\textbf{A fine-tuning claim deserves a more precise treatment than v1 gave it.} v1 asserted that fine-tuning "primarily buys" style and fluency rather than factual accuracy --- too categorical a claim. The more defensible, narrower proposition: ordinary fine-tuning is often an unreliable way to inject large amounts of \emph{new} factual knowledge into a model, and retrieval frequently performs better for that specific purpose. But this depends heavily on what kind of fine-tuning is being discussed --- supervised task adaptation, continued pretraining, preference tuning, and deliberate factual-knowledge injection are different techniques with different track records, and fine-tuning \emph{can} improve both factuality and task accuracy when the training objective and data are specifically designed to do that. A tooling vendor's claim ("we trained it on 500,000 filings") should be evaluated on the evidentiary question --- does this measurably improve accuracy on real tasks, or only confidence --- not dismissed by a blanket rule about what fine-tuning generally does.
